\chapter{State of the Art}
\label{ch:soa}

This chapter surveys the work that motivates, contextualises, and positions
the system developed in this thesis. It is organised into five areas:
commercial AR interior design applications and their limitations
(Section~\ref{sec:soa:commercial}); indoor scene understanding datasets and
methods that establish the semantic reconstruction context
(Section~\ref{sec:soa:indoor}); dense SLAM and volumetric reconstruction
systems that form the technological lineage of depth-based room capture
(Section~\ref{sec:soa:slam}); mobile AR usability evaluation methodology
(Section~\ref{sec:soa:usability}); and the evolution of novel-view synthesis
from neural radiance fields to 3D Gaussian Splatting
(Section~\ref{sec:soa:nvs}). Section~\ref{sec:soa:gap} synthesises the
research gap that this thesis addresses.

% ------------------------------------------------------------------
\section{Commercial AR Interior Design Applications}
\label{sec:soa:commercial}
% ------------------------------------------------------------------

A number of consumer-facing augmented reality tools allow users to preview
furniture in their physical environment before purchase. The most widely
deployed example is IKEA~Place, released in 2017 as one of the first
ARKit-based shopping applications. It uses ARKit's plane detection to
identify the floor surface and anchors a scaled virtual product model to that
plane, allowing the user to walk around a composited sofa or bookcase in the
live camera feed. Similar mechanics underpin the AR viewers embedded in the
Wayfair, Houzz, and Amazon applications. These tools have demonstrably
reduced purchase uncertainty for a subset of users, but they share a
structural limitation: the room is not modelled. Plane detection locates one
or more horizontal surfaces; it does not produce walls, identify windows or
doors, record the positions of existing furniture, or generate any artefact
that persists beyond the session. Each product preview starts from scratch,
and the user must be physically inside the room to perform it.

K{\'a}n~et~al.~\cite{kan2021ar} address a complementary limitation: even
when an AR tool allows placement, it provides no layout intelligence. Their
system constructs a hierarchical tree of procedural rules encoding spatial
relations and style constraints, then uses a recommender engine to propose
furniture arrangements that satisfy those rules; the AR component overlays
the suggested layout on the live camera feed. This work goes substantially
further than product-viewer apps in terms of automated design support.
However, it does not start from a captured parametric model of the existing
room: the constraint solver operates in an abstract space whose connection to
metric room geometry depends on the quality of transient plane detection, and
no mechanism for off-site editing is described.

Research into AR interior workflows with semantic scene understanding has
explored using depth-sensing cameras and learned 3D object detectors to make
virtual object grounding more robust~\cite{Geo2023}. Pugh~et~al.~\cite{Geo2023}
demonstrate that synthetic indoor datasets with photorealistic rendering and
rich semantic labels can substantially improve indoor scene understanding for
AR applications, an approach complementary to the real-scan pipeline used
in this thesis.

The gap that all these systems leave open is the absence of a single
application that (i)~captures the room as a structured, metric, persistent
model; (ii)~allows editing of that model independently of being on-site;
and (iii)~generates or imports three-dimensional assets through an on-device
photogrammetric pipeline. This thesis addresses that gap.

% ------------------------------------------------------------------
\section{Indoor Scene Understanding and Semantic Datasets}
\label{sec:soa:indoor}
% ------------------------------------------------------------------

The feasibility of semantic room capture on consumer hardware is grounded in
a body of indoor scene understanding research that established both datasets
and benchmarks.

ScanNet~\cite{dai2017scannet}, introduced by Dai~et~al., comprises 1\,513
RGB-D scans of real indoor environments reconstructed offline and annotated
with per-vertex semantic labels across 20 object categories. ScanNet
demonstrated that rich 3D semantics are achievable at scale with commodity
depth sensors and enabled a generation of learning-based methods for 3D
object detection, semantic segmentation, and scene completion. Crucially, it
showed that labelled 3D structure, knowing not only where a surface is but
also \emph{what} it is, dramatically improves downstream task performance
compared to raw geometry alone.

Matterport3D~\cite{CAS+17}, introduced by Chang~et~al., extended this to
building scale: 90 buildings comprising 10\,800 panoramic RGB-D views and
surface reconstructions, with instance-level semantic labels for 40
categories. The larger scale and higher-resolution surface models of
Matterport3D enabled research on room-level layout understanding, floor plan
estimation, and cross-room navigation, tasks relevant to applications that
must understand not just a single room but a home or office as a whole.

These datasets established that the semantic granularity required for useful
interior design support, distinguishing structural surfaces from furniture,
classifying furniture by category, estimating bounding dimensions, is
achievable in principle on real hardware at real scale. Apple's RoomPlan
API~\cite{apple2024roomplan2} brings a consumer subset of this capability to
supported devices, outputting a \texttt{CapturedRoom} whose typed surfaces
and object instances reflect the same conceptual taxonomy that ScanNet and
Matterport3D made trainable. The key difference is the deployment context:
whereas ScanNet's pipeline required offline reconstruction and manual
annotation over many hours, RoomPlan runs entirely on-device in real time
during an interactive scanning session of one to three minutes.

% ------------------------------------------------------------------
\section{Dense SLAM and Volumetric Reconstruction}
\label{sec:soa:slam}
% ------------------------------------------------------------------

The technological lineage of depth-based room capture runs through a series
of dense simultaneous localisation and mapping (SLAM) systems that
progressively made real-time volumetric reconstruction practical on
increasingly constrained hardware.

\textbf{KinectFusion}~\cite{newcombe2011kinectfusion}, introduced by
Newcombe~et~al.\ at ISMAR~2011, was the first system to demonstrate
real-time dense surface reconstruction using a consumer depth sensor. It
represents the scene as a truncated signed distance function (TSDF) stored
in a fixed-resolution voxel grid: each incoming depth frame updates the TSDF
along viewing rays using a weighted running average, and a tracking step
registers each frame to the current model using Iterative Closest Point
(ICP) on projective depth residuals. KinectFusion established the TSDF as
the standard representation for dense RGB-D SLAM and showed that noise in
structured-light depth measurements averages down effectively over many
frames, producing coherent surfaces from individually noisy observations.
Its limitation was scale: a fixed voxel volume of, typically, $512^3$ cells
cannot represent scenes larger than a few cubic metres without reducing
resolution or exhausting GPU memory.

\textbf{ElasticFusion}~\cite{WLS+15}, by Whelan~et~al., addressed the loop
closure problem that arises in room-scale capture: when the sensor returns
to a previously observed area after accumulating pose drift, the map
develops misalignment seams. Rather than correcting these with a pose graph
and rigid re-projection, ElasticFusion applies a non-rigid surface
deformation to bring the entire map into global consistency when a loop is
detected. The result is a dense, coloured surface reconstruction of
room-scale environments without the visible discontinuities characteristic
of pose-graph approaches.

\textbf{BundleFusion}~\cite{DAN+17}, by Dai~et~al., pushed the state of
the art further by combining dense geometric alignment with sparse
feature-based re-localisation in a unified online optimisation. BundleFusion
operates in real time and produces globally consistent reconstructions of
full apartment-scale scenes through on-the-fly surface
reintegration, re-fusing surface observations whenever a new global
alignment is computed. It was used as the reconstruction backend for the
ScanNet dataset~\cite{dai2017scannet}, establishing a direct link between
the volumetric reconstruction literature and the semantic dataset research
reviewed in Section~\ref{sec:soa:indoor}.

\textbf{ORB-SLAM3}~\cite{CEG+21}, by Campos~et~al., represents the current
state of the art in open-source visual and visual-inertial SLAM. Although
primarily a sparse mapping system that builds a landmark-based map rather
than a dense surface, ORB-SLAM3 is relevant here because ARKit's
Visual-Inertial Odometry (VIO), which underpins both Object Capture and
RoomPlan, uses structurally similar principles: a tightly coupled
camera--IMU estimator maintaining a sliding-window bundle adjustment over
recent keyframes. Understanding this VIO tracking layer clarifies both the
accuracy characteristics and the failure modes of RoomPlan scans,
particularly in textureless environments or during rapid device motion.

The progression from KinectFusion~\cite{newcombe2011kinectfusion} through
ElasticFusion~\cite{WLS+15} and BundleFusion~\cite{DAN+17} to ARKit-based
capture represents a reduction of hardware requirements from a tethered GPU
workstation to an iPhone held in one hand, at the cost of some geometric
resolution and flexibility. RoomPlan adds semantic interpretation on top of
this depth fusion foundation, trading raw surface detail for the structured,
editable output that makes furniture replacement possible.

% ------------------------------------------------------------------
\section{Mobile AR Usability Evaluation}
\label{sec:soa:usability}
% ------------------------------------------------------------------

The design decisions in this thesis, in particular the choice to provide a
non-AR, persistence-first editor alongside live AR previewing, are informed
by empirical findings on how users interact with mobile AR interior
applications.

Kandil~et~al.~\cite{kandil2021ux} conducted a comparative evaluation of
five commercial AR interior design applications using two complementary
instruments: the System Usability Scale (SUS) and the Handheld Augmented
Reality Usability Scale (HARUS). SUS is a widely used 10-item Likert
questionnaire for measuring overall perceived usability; a score above 70 is
conventionally considered acceptable and above 85 is
excellent~\cite{kandil2021ux}. HARUS extends SUS with items specific to
handheld AR, addressing factors such as physical fatigue from sustained
device holding, accuracy of virtual object placement relative to real
geometry, and the cognitive load of simultaneously navigating real and
virtual spaces.

Kandil~et~al.\ found that across the five tested applications, the
lowest-scoring dimensions were \emph{manipulation precision}, how
accurately a user could position and orient a virtual object, and
\emph{environmental awareness}, how well the AR composite preserved the
user's spatial understanding of the real room. Both failure modes are
directly addressed by the design of this thesis: manipulation precision
improves when editing is performed in a stable 3D scene with a free-moving
camera rather than in a live AR feed where the device must be held up and
pointed; environmental awareness improves when virtual objects are grounded
to a captured, metric room model rather than to a transiently detected floor
plane. The study also noted that persistent project state, the ability to
save a staging session and return to it later, was mentioned by participants
as a desired feature absent from all evaluated applications, motivating
Requirement~R-2 of this thesis directly.

Although this thesis does not itself conduct a formal SUS/HARUS evaluation,
the usability criteria identified by Kandil~et~al.~\cite{kandil2021ux}
inform the non-functional requirements specified in
Chapter~\ref{ch:requirements}, and the methodology described
in~\cite{kandil2021ux} provides a ready-made framework for a future user
study.

% ------------------------------------------------------------------
\section{Novel View Synthesis: From NeRF to 3D Gaussian Splatting}
\label{sec:soa:nvs}
% ------------------------------------------------------------------

Novel-view synthesis, rendering a scene from a viewpoint not present in the
training images, underwent a step change in 2020 with the introduction of
Neural Radiance Fields (NeRF) by Mildenhall~et~al.~\cite{mildenhall2020nerf}.
NeRF represents a scene as a continuous volumetric function
$F_\Theta : (\mathbf{x}, \mathbf{d}) \mapsto (\mathbf{c}, \sigma)$
approximated by a multi-layer perceptron (MLP), mapping a 3D position
$\mathbf{x}$ and viewing direction $\mathbf{d}$ to colour $\mathbf{c}$ and
volume density $\sigma$. Pixel colours are computed by numerical integration
of the classical volume rendering equation along camera rays. Trained on a
set of posed images, NeRF achieves previously unattained view synthesis
quality, capturing fine view-dependent effects such as specular highlights
and semi-transparency. The two limitations of the original formulation are
training time (tens of hours per scene) and rendering speed (several seconds
per frame), both of which prevent interactive use.

Instant Neural Graphics Primitives (Instant-NGP)~\cite{MESK22}, introduced
by M\"{u}ller~et~al., reduced training to seconds by replacing the MLP's
positional encoding with a multiresolution hash grid that maps 3D positions
to compact learned feature vectors. The MLP itself becomes shallow, the
hash grid captures most of the spatial variation, and the total parameter
count drops dramatically. Instant-NGP brought NeRF-quality novel-view
synthesis within reach for interactive applications, though rendering
remained ray-marching based and therefore not real-time on commodity
hardware.

Mip-NeRF~360~\cite{BMV+22}, by Barron~et~al., addressed quality limitations
in unbounded indoor and outdoor scenes. Standard NeRF uses a point sample
along each ray, which causes aliasing when the same scene element is
observed at very different scales; Mip-NeRF~360 uses cone frustum samples
and integrates the encoding over the cone, substantially reducing aliasing
artefacts. It also introduces a parameterisation for unbounded scene extent
using a contracted coordinate space. Mip-NeRF~360 currently defines one of
the quality upper bounds against which subsequent methods, including 3DGS,
are evaluated.

\textbf{3D Gaussian Splatting} (3DGS)~\cite{kerbl20233dgs}, introduced by
Kerbl~et~al.\ at SIGGRAPH~2023, replaces the neural implicit representation
with an explicit cloud of anisotropic Gaussian primitives. Each primitive
stores a mean $\boldsymbol{\mu} \in \mathbb{R}^3$, a covariance $\Sigma$
decomposed as $\Sigma = RSS^\top R^\top$ (rotation $R$, diagonal scale $S$),
an opacity $\alpha$, and view-dependent colour encoded as low-degree
spherical harmonic (SH) coefficients. SH provide a compact, differentiable
approximation to view-dependent radiance: degree-0 SH captures diffuse
colour; higher degrees capture glossy and specular effects. Rendering
projects each Gaussian onto the image plane via the Jacobian of the
perspective mapping, sorts the resulting footprints by depth, and composites
them with the standard alpha-blending operator. Because this is a
rasterisation operation rather than ray marching, Kerbl~et~al.\ report
rendering speeds exceeding 100\,FPS at 1080p on a desktop GPU while
matching or exceeding Mip-NeRF~360 quality on standard
benchmarks~\cite{kerbl20233dgs}. Training initialises from a COLMAP sparse
point cloud, linking 3DGS directly to the Structure-from-Motion, and optimisation alternates
gradient steps with adaptive density control: large Gaussians are split to
capture fine detail, and nearly-transparent primitives are pruned to control
memory.

For mobile deployment, PocketGS~\cite{Guo26} demonstrates on-device 3DGS
training on iPhone hardware, showing that the Gaussian representation is
approaching the constraints of the platform used in this thesis. In the
current application, training is off-device; the Metal-based renderer
described in Chapter~\ref{ch:implementation} loads a pretrained
\texttt{.ply} file. The architecture is designed so that if on-device
training becomes practical, the import step can be replaced without
modifying the renderer.

% ------------------------------------------------------------------
\section{Summary and Research Gap}
\label{sec:soa:gap}
% ------------------------------------------------------------------

Table~\ref{tab:soa:comparison} positions this thesis against the four
categories of related work surveyed above across dimensions that are
directly relevant to practical interior design support on mobile hardware.

\begin{table}[ht]
\centering
\small
\caption{Positioning of this thesis relative to representative related
systems. \checkmark~=~supported; $\circ$~=~partial or session-only;
$\times$~=~not supported.}
\label{tab:soa:comparison}
\begin{tabular}{lcccc}
\toprule
\textbf{System} &
\makecell{\textbf{Metric}\\\textbf{room model}} &
\textbf{Persistence} &
\makecell{\textbf{On-device}\\\textbf{assets}} &
\makecell{\textbf{Splat}\\\textbf{viewing}} \\
\midrule
IKEA Place / Wayfair AR                 & $\times$   & $\times$   & $\times$   & $\times$ \\
K{\'a}n~et~al.~\cite{kan2021ar}         & $\circ$    & $\times$   & $\times$   & $\times$ \\
ScanNet pipeline~\cite{dai2017scannet}  & \checkmark & \checkmark & $\times$   & $\times$ \\
3DGS~\cite{kerbl20233dgs} (standalone) & $\times$   & \checkmark & $\times$   & \checkmark \\
\textbf{This thesis}                    & \checkmark & \checkmark & \checkmark & \checkmark \\
\bottomrule
\end{tabular}
\end{table}


The survey reveals a consistent gap: existing consumer AR tools either lack
a structured room model entirely, or possess one but offer no interactive
editing and no on-device asset authoring. Research systems such as the
ScanNet reconstruction pipeline~\cite{dai2017scannet} produce high-quality
semantic models but require offline processing pipelines unsuitable for
end-user workflows. Procedural layout systems such as
K{\'a}n~et~al.~\cite{kan2021ar} add design intelligence but not persistent
capture. 3DGS toolchains~\cite{kerbl20233dgs} produce high-fidelity scene
representations but no parametric structure suitable for furniture
replacement.

This thesis occupies the intersection: Apple's on-device
APIs~\cite{apple2024roomplan2} provide a metric, typed room model;
RealityKit provides an editable, persistent scene graph; Object Capture
provides on-device photogrammetric asset authoring; and a Metal-based splat
renderer provides a forward-looking viewing path for the Gaussian
representation. The usability evaluation methodology of
Kandil~et~al.~\cite{kandil2021ux} provides the framework for validating
the result in a future formal study.
